% ============================================================
%  CHAPTER 1 — Introduction
%  Logical chain: Problem → Gap → Proposed Solution
%  Status: Not started
% ============================================================

\chapter{Introduction}
\label{chap:introduction}

% ─── 1.1 MOTIVATION ─────────────────────────────────────────
\section{Motivation}
\label{sec:motivation}

% TODO: Motivate WBC for humanoid robots.
% Argue: Why does humanoid whole-body coordination matter?
% Reference: service robotics, industrial assistance contexts.
% Cite key humanoid platforms (Atlas, TALOS, Unitree H1/G1).

\lipsum[2]

\begin{criticalinsight}
  % TODO: Identify the core tension that makes this problem hard.
  % e.g., The gap between kinematic redundancy and real-time constraint satisfaction.
\end{criticalinsight}

% ─── 1.2 PROBLEM STATEMENT ──────────────────────────────────
\section{Problem Statement}
\label{sec:problem}

% TODO: Formalize the research problem.
% Be precise: what specific limitation in existing WBC frameworks motivates this work?
% Anchor to: absence of a standard framework coupling planning + QP WBC for
% simultaneous loco-manipulation (master_memory §7.2 — Methodological Gap).

\lipsum[3]

\begin{criticalinsight}
  % TODO: State the falsifiable core claim.
\end{criticalinsight}

% ─── 1.3 RESEARCH QUESTION ──────────────────────────────────
\section{Research Question and Objectives}
\label{sec:research-question}

% Canonical research question (master_memory §1.2):
\begin{quote}
  \textit{``How can a \WBC{} framework be designed and validated in simulation
  to enable a humanoid robot to perform stable, multi-contact locomotion and
  manipulation tasks under dynamic constraints?''}
\end{quote}

% TODO: Refine once formally agreed with supervisor (PQ4).

\noindent The following sub-questions structure the investigation:
\begin{enumerate}[label=\textbf{SQ\arabic*}]
  \item What are the key limitations of existing WBC formulations (task-priority,
        QP-based, MPC) when applied to full-body humanoid coordination?
  \item How can contact dynamics be modeled and integrated in a simulation
        environment to faithfully replicate hardware behavior?
  \item What metrics and scenarios constitute a rigorous simulation-based
        validation of a WBC framework for humanoids?
\end{enumerate}

% ─── 1.4 HYPOTHESES ─────────────────────────────────────────
\section{Hypotheses}
\label{sec:hypotheses}

\begin{hypothesis}{H1}
  A hierarchical \QP{}-based \WBC{} formulation can achieve stable whole-body
  coordination across locomotion and manipulation tasks simultaneously.
\end{hypothesis}

\begin{hypothesis}{H2}
  Simulation-based validation (Isaac Sim / MuJoCo) is sufficient to demonstrate
  proof-of-concept for the proposed WBC framework before hardware deployment.
\end{hypothesis}

\begin{hypothesis}{H3}
  Explicit contact modeling and task-priority hierarchies outperform naive
  torque control in multi-contact transitions.
\end{hypothesis}

% ─── 1.5 SCOPE AND LIMITATIONS ──────────────────────────────
\section{Scope and Limitations}
\label{sec:scope}

% TODO: Summarise the in-scope / out-of-scope table (master_memory §1.5).
% Explicitly exclude: hardware experiments, wheeled robots, end-to-end RL as
% primary contribution (master_memory §11 — Rejected Approaches).

\lipsum[4]

% ─── 1.6 CONTRIBUTIONS ──────────────────────────────────────
\section{Contributions}
\label{sec:contributions}

% TODO: Clearly delineate novelty vs. Crocoddyl / existing WBC stacks (PQ4).
% Candidate contributions:
%  (1) Novel hierarchical coupling of centroidal MPC + QP execution for
%      simultaneous loco-manipulation.
%  (2) Reproducible simulation-based evaluation protocol addressing the
%      evaluation gap (master_memory §7.3).
%  (3) Comparative benchmark of WBC, task-priority, and RL baselines on
%      standardized S1–S5 scenarios.

\lipsum[5]

% ─── 1.7 THESIS STRUCTURE ───────────────────────────────────
\section{Thesis Structure}
\label{sec:structure}

The remainder of this thesis is organized as follows.
\Cref{chap:sota} reviews the state of the art in whole-body control for humanoid
robots, covering model-based methods, centroidal MPC, and learning-based baselines.
\Cref{chap:methodology} formalizes the proposed WBC framework and justifies the
design choices.
\Cref{chap:simulation} describes the simulation environment, robot model, and
software stack.
\Cref{chap:experiments} presents quantitative results across evaluation scenarios
S1 through S5.
\Cref{chap:discussion} interprets the results and discusses limitations.
\Cref{chap:conclusion} summarizes contributions and outlines future research
directions.
